\begin{afterword}
\summaryhead

A short story. Jeffreyssai tells his class they must aspire to do better than Einstein, whose ten years on General Relativity were ``too slow.'' He sets five students the correct theory of quantum gravity, in one month, as a test of speed, and leaves the room.

Styrlyn takes the lead. Each student writes down angles alone, then gives one, with criticism held back. Taji says to assume Eld science's avenues were blind alleys and to avoid string theory; Hiriwa, to get rid of the infinities; Yin, citing a Lojban inscription, to eliminate time from the equations; Styrlyn, to find a spatially local form of quantum physics. Brennan offers ``the obvious'': why is the energy of the quantum Hamiltonian necessarily the energy that curves spacetime? Hiriwa rephrases it as a search for a view in which the two are one. After thirty-seven seconds of silence someone says, ``I have an idea\ldots'' and the story ends.

\responsehead

The story is fiction and is judged for what it shows. It stages two earlier posts. Jeffreyssai's opening is the thesis of ``Einstein's Superpowers,'' posted the day before, and the one-month deadline comes from ``The Failures of Eld Science.'' His charges against Einstein repeat the author's own, and one of them, that Einstein ``reasoned cleanly'' on General Relativity, passes over the two years he spent on a wrong theory.

What it shows well is the method. The students write ideas alone before combining them and hold criticism back, an order that research on brainstorming supports. Their angles address real difficulties: the infinities of quantized gravity and the different roles of time in quantum theory and relativity are standard problems in the field. Brennan's question is modeled on Einstein's step of making two masses one quantity. Yin's worry that it is ``too obvious'' has some basis. In 1964 Weinberg derived the equality of gravitational and inertial mass from quantum theory, a proof of the kind Taji proposes and Hiriwa sets aside.

The framing leaves one question open. ``The Failures of Eld Science'' asked whether a problem could be solved in a month in a case where the evidence is already in hand. Quantum gravity has, as yet, no experimental tests, and the story does not say whether the students have enough to work with.

Like ``The Ritual,'' the story stops before any result, so the reader sees the method and the stakes but not whether the speed Jeffreyssai demands can be had.

In short: a brisk scene that stages the previous day's thesis, shows a sound way to gather ideas, and gives its students angles that address real problems of quantum gravity; like ``The Ritual,'' it stops before any result.
\end{afterword}
